\section{Feasibility certificates beyond the second moment}
\label{sec:convex-feasibility}

The scalar second-moment test is one member of a larger family of
path-level restrictions.  These restrictions can reject a path before a
nonlinear correction search and give positive lower bounds on the best
possible acceleration residual.  The underlying second-derivative chain
rule is standard transport calculus; here it yields certificates that
strictly strengthen the second-moment restriction.

For fixed $t$, define the dual norm
\begin{equation}
 \varepsilon_t^{\rm acc}
 =\sup_{\eta\in C_c^\infty,\,\|\nabla\eta\|_{L^2(p_t)}>0}
 \frac{\left|\int \nabla\eta\cdot\acc_t\,\dd p_t\right|}
 {\|\nabla\eta\|_{L^2(p_t)}}.
 \label{eq:acc-dual-norm}
\end{equation}
It equals the norm of the weighted-gradient projection of $\acc_t$ and is
no larger than $\|\acc_t-w\|_{L^2(p_t)}$ for every $w$ of zero weighted
divergence.  This identity alone asserts no regularity of that projection.

\begin{theorem}[Convex-observable and weighted-energy certificates]
\label{thm:convex-feasibility}
Assume the regularity and moment hypotheses of
Theorem~\ref{thm:distribution_order}, except for the acceleration condition.
For $\phi\in C^2(\R^d)$ with bounded Hessian, set
$F_\phi(t)=\int\phi\,\dd p_t$.  Then
\begin{equation}
 F_\phi''(t)=\int u_t^\top D^2\phi\,u_t\,\dd p_t
       +\int\nabla\phi\cdot\acc_t\,\dd p_t.
 \label{eq:convex-chain-rule}
\end{equation}
If the acceleration condition holds and $\phi$ is convex, $F_\phi$ is
convex in time.  More generally, let $K_\phi(t)$ be the infimum of
$\int v^\top D^2\phi\,v\,\dd p_t$ over finite-energy instantaneous velocities
satisfying $\partial_t p_t+\nabla\cdot(p_tv)=0$.  For convex $\phi$ and
$\|\nabla\phi\|_{L^2(p_t)}>0$,
\begin{equation}
 \varepsilon_t^{\rm acc}\ge
 \frac{[K_\phi(t)-F_\phi''(t)]_+}
 {\|\nabla\phi\|_{L^2(p_t)}}
 \ge\frac{[-F_\phi''(t)]_+}{\|\nabla\phi\|_{L^2(p_t)}}.
 \label{eq:convex-residual-floor}
\end{equation}
Unbounded Hessians are permitted when the chain rule is dominated and
$\nabla\phi$ belongs to the weighted-gradient closure, including
polynomial tests under the corresponding finite moments.
\end{theorem}

Thus a feasible path satisfies
$p_{(1-\lambda)s+\lambda t}\preceq_{\rm cx}(1-\lambda)p_s+\lambda p_t$:
each convex observable at an intermediate time is bounded by its endpoint
mixture average.  This is necessary, not sufficient.  Every straight
stochastic interpolation $\Law((1-t)X_0+tX_1)$ already passes this test.
Under weighted Poincare and uniform ellipticity assumptions, $K_\phi$
is obtained by a linear weighted elliptic problem.
Appendix~\ref{app:convex-feasibility} states these hypotheses and gives
the characterization and all proofs.

\begin{proposition}[Covariance constraints]
\label{prop:matrix-feasibility}
Under the acceleration condition and the preceding hypotheses, let
$m(t)=m(0)+t\xi$ and $C(t)=\operatorname{Cov}(X_t)>0$.  There exists $B_t$
such that
\begin{equation}
 C'=B_t+B_t^\top,\qquad
 \begin{pmatrix}C&B_t^\top\\ B_t&C''/2\end{pmatrix}\succeq0.
 \label{eq:covariance-lmi}
\end{equation}
In particular $C''\succeq0$, and every directional standard deviation
$\sqrt{e^\top C(t)e}$ is convex.
\end{proposition}

The block-matrix condition is a semidefinite necessary test using only
the prescribed covariance and its first two derivatives.  It can fail
even when the scalar second-moment test passes.  The following example
passes both tests but is excluded by a fourth-order convex observable.

\begin{proposition}[A residual floor beyond covariance tests]
\label{prop:quartic-obstruction}
Let $Z\sim\mathcal N(0,I_2)$ and $0\le t\le1/2$.  Define
\begin{equation}
 \begin{split}
 \kappa(r)&=\frac{-1+5e^{-r^2/2}}4,\qquad b(z)=\kappa(\|z\|)z,\\
 G_t(z)&=z+\frac{t^2}{2}b(z),\qquad p_t=(G_t)_\#\Law(Z).
 \end{split}
 \label{eq:quartic-obstruction-path}
\end{equation}
This path has smooth positive densities and a smooth globally Lipschitz
minimum-action generator $v_t(G_tz)=t b(z)$.  It passes the scalar
second-moment and covariance tests throughout the interval.  Nevertheless,
no generator with the stated regularity can satisfy the acceleration
condition.  Every such generator obeys
\begin{equation}
 \varepsilon_0^{\rm acc}\ge\frac{\sqrt3}{16}.
 \label{eq:quartic-floor}
\end{equation}
\end{proposition}

The separation is explicit:
$M''-2K_{\min}=1/4+23t^2/144>0$, and
$C(t)=(1+t^2/16+23t^4/1152)I_2$ satisfies
\eqref{eq:covariance-lmi}.  However,
$(\E\|X_t\|^4)''<0$ on the entire interval, with value $-3$ at zero.
Since $\|\nabla\|x\|^4\|_{L^2(p_0)}=16\sqrt3$,
\eqref{eq:convex-residual-floor} gives the displayed floor.  The obstruction
applies to every admissible generator, independently of its
parameterization or optimization method.  In particular no correction
family on this path can achieve
$\sup_t\varepsilon_t^{\rm acc}=O(H)$ as $H\downarrow0$.
By Corollary~\ref{cor:sharp-local-residual}, the same obstruction yields
an actual local error floor for every such generator, with
$E^u_{0,h}(x)=x+hu(0,x)$:
\[
 \lim_{h\downarrow0}h^{-2}W_2((E^u_{0,h})_\#p_0,p_h)
 \ge\frac{\sqrt3}{32}.
\]
This local conclusion does not exclude accidental cancellation of
terminal integration error.
